\documentclass[11pt, a4paper]{article}
\usepackage[utf8]{inputenc}
\usepackage[T1]{fontenc}
\usepackage[french]{babel}
\usepackage{lmodern}
\usepackage{geometry}
\geometry{margin=2.5cm}
\usepackage{hyperref}
\hypersetup{
    colorlinks=true,
    linkcolor=blue,
    urlcolor=blue,
    bookmarksnumbered=true,
    pdftitle={Manuel d'Utilisation Pangol1},
    pdfauthor={Pangol1 Team}
}
\usepackage{enumitem}
\usepackage{graphicx}
\usepackage{xcolor}
\usepackage{float}
\usepackage{tabularx}
\usepackage{booktabs}

\title{Manuel d'Utilisation : Pangol1}
\author{}
\date{}

\addto\captionsfrench{\renewcommand{\contentsname}{Sommaire}}

\setlength{\parindent}{0pt}
\setlength{\parskip}{0.35em}
\setlist{topsep=0.2em, itemsep=0.2em}

\newcommand{\placeholderfigure}[3]{%
\begin{figure}[htbp]
    \centering
    \fbox{\parbox[c][0.32\textheight][c]{0.9\textwidth}{\centering
        \textit{#1}\\[0.4em]
        \texttt{#2}}}
    \caption{#1}
    \label{fig:#3}
\end{figure}
}

\begin{document}

\maketitle

\section*{À propos de ce document}
\addcontentsline{toc}{section}{À propos de ce document}
Ce manuel décrit l'interface graphique de \textbf{Pangol1}, les actions principales et les réglages disponibles. Il a été conçu pour accueillir des captures d'écran et des schémas supplémentaires dans les sections clés.

\vspace{2ex}
\noindent
{\small \textit{Dernière mise à jour : \today}}

\newpage

\tableofcontents

\newpage

\section{Introduction}

\subsection{Qu'est-ce que Pangol1 ?}
\textbf{Pangol1} est une application graphique permettant d'explorer des données tabulaires et de générer des rapports visuels (graphiques, images et texte) exportables.

\subsection{Formats pris en charge}
\begin{itemize}
    \item Fichiers \textbf{CSV}, \textbf{TSV} et \textbf{TXT}
    \item Archives \textbf{ZIP} (extraction puis chargement)
    \item Fichiers \textbf{Parquet}
    \item Chargement depuis une \textbf{URL}
\end{itemize}

\subsection{Fonctionnalités principales}
\begin{itemize}
    \item Chargement multi-fichiers et gestion des tables récentes
    \item Table interactive (pagination, tri, filtres, masquage et renommage)
    \item Statistiques détaillées par colonne et indicateurs rapides en barre d'état
    \item Création d'un rapport (Notebook) avec graphiques, images et texte
    \item Export du rapport en HTML, PDF ou Markdown
    \item Thèmes, langues et réglages d'interface personnalisables
\end{itemize}

\subsection{Conventions de lecture}
\begin{itemize}
    \item Les menus sont indiqués en gras, par exemple : \textbf{Fichier → Ouvrir}
    \item Les raccourcis clavier sont indiqués en style \texttt{Ctrl+O}
\end{itemize}

\section{Démarrage rapide}

\subsection{Lancer l'application}
Pangol1 fonctionne sur Windows, Linux et macOS.
\begin{enumerate}
    \item Lancez \textbf{Pangol1.jar} ou exécutez \texttt{java -jar Pangol1.jar}
    \item En développement, exécutez la classe principale \textbf{Pangol1}
    \item L'interface principale s'affiche après l'écran de chargement
\end{enumerate}

\subsection{Première utilisation}
\begin{enumerate}
    \item Ouvrez un fichier via \textbf{Fichier → Ouvrir}
    \item Les données s'affichent dans la table
    \item Utilisez le panneau de gauche pour créer un graphique ou un rapport
\end{enumerate}

\section{Vue générale de l'interface}

\placeholderfigure{Vue générale de l'interface}{images/ui-overview.png}{ui-overview}

\subsection{Zones principales}
\begin{itemize}
    \item \textbf{Barre de menus} : accès aux fichiers, affichages et aides
    \item \textbf{Panneau de rapport (gauche)} : création et édition du Notebook
    \item \textbf{Table de données (droite)} : exploration et filtrage des données
    \item \textbf{Barre d'état (bas)} : informations de statut et statistiques rapides
\end{itemize}

\subsection{Séparateurs (split)}
Les panneaux sont séparés par des barres réglables. L'orientation principale et celle du Notebook peuvent être modifiées via \textbf{Affichage → Split}.

\section{Barre de menus}

\subsection{Menu Fichier}
\begin{itemize}
    \item \textbf{Ouvrir} : charger un ou plusieurs fichiers (CSV, TSV, TXT, ZIP, Parquet)
    \item \textbf{Ouvrir récent} : accéder aux tables récemment ouvertes
    \item \textbf{Ouvrir une URL} : charger un fichier distant et choisir un dossier de destination
    \item \textbf{Paramètres} : ouvrir la fenêtre de réglages
    \item \textbf{Quitter} : fermer l'application
\end{itemize}

\subsection{Menu Affichage}
\begin{itemize}
    \item \textbf{Panneaux} : afficher ou masquer la table, le Notebook et le panneau de réglages du Notebook
    \item \textbf{Split} : orientation principale, orientation du Notebook, réinitialisation (principal, Notebook, complet)
\end{itemize}

\subsection{Menu Aide}
\begin{itemize}
    \item \textbf{Langues} : basculer l'interface (Français, Anglais, Espagnol, Arabe, Portugais, Russe, Japonais, Indonésien, Turc, Chinois simplifié)
    \item \textbf{Documentation} : ouvrir la documentation (si configurée)
    \item \textbf{À propos} : version, licence et équipe
\end{itemize}

\section{Chargement et gestion des données}

\subsection{Ouvrir un fichier}
\begin{enumerate}
    \item \textbf{Fichier → Ouvrir}
    \item Sélectionnez un ou plusieurs fichiers
    \item Pour les \textbf{ZIP}, choisissez un dossier de destination pour l'extraction
    \item Si plusieurs tables sont chargées, elles apparaissent dans la liste des tables
\end{enumerate}

\subsection{Ouvrir une URL}
\begin{enumerate}
    \item \textbf{Fichier → Ouvrir une URL}
    \item Saisissez une URL valide
    \item Choisissez un dossier où stocker le fichier téléchargé
\end{enumerate}

\subsection{Écran vide}
Lorsque aucune table récente n'est disponible, un bouton \textbf{Charger} permet d'ouvrir un fichier.

\subsection{Liste des tables récentes}
Lorsque aucune table n'est ouverte, Pangol1 affiche une liste des tables chargées. Un double-clic ouvre une table. Un clic droit permet :
\begin{itemize}
    \item \textbf{Renommer} l'alias d'une table
    \item \textbf{Ouvrir l'emplacement} du fichier
    \item \textbf{Retirer des récents}
    \item \textbf{Supprimer du disque}
\end{itemize}

\subsection{Fermer une table}
Utilisez le bouton de fermeture dans la barre d'outils de la table. La vue bascule sur la liste ou l'écran vide.

\section{Table de données}

\placeholderfigure{Vue table de données}{images/table-view.png}{table-view}

\subsection{Pagination}
\begin{itemize}
    \item Navigation via \textbf{First, Previous, Next, Last}
    \item Taille de page configurable : 500, 1000, 5000, 10000 lignes
    \item L'option de taille n'apparaît que pour les tables volumineuses
\end{itemize}

\subsection{Barre d'outils}
\begin{itemize}
    \item \textbf{Filtres avancés} : ouvre la fenêtre de filtrage multi-conditions
    \item \textbf{Effacer les filtres}
    \item \textbf{Afficher toutes les colonnes}
    \item \textbf{Masquer les colonnes vides}
    \item \textbf{Fermer la table}
\end{itemize}

\subsection{En-têtes et sélection}
\begin{itemize}
    \item Clic gauche sur un en-tête : sélection de la colonne
    \item Clic droit sur un en-tête : menu contextuel complet
    \item Un pictogramme indique qu'un filtre est actif sur une colonne
\end{itemize}

\subsection{Menu contextuel d'une colonne}
\begin{itemize}
    \item \textbf{Tri} : ascendant ou descendant
    \item \textbf{Filtrage} : Top N et Bottom N (longueur, valeur ou date selon le type), plage de valeurs, vides, non vides
    \item \textbf{Statistiques} : résumé, taux de valeurs nulles, cardinalité, IQR, asymétrie, variation, corrélation
    \item \textbf{Renommer} et \textbf{Masquer} la colonne
\end{itemize}

\subsection{Statistiques rapides}
Lorsque vous sélectionnez une colonne, la barre d'état affiche \textbf{min}, \textbf{max}, \textbf{moyenne} et \textbf{somme} si disponibles.

\section{Filtrage avancé}

\subsection{Créer un filtre}
\begin{enumerate}
    \item Ouvrez la fenêtre via \textbf{Filtres avancés}
    \item Choisissez une colonne, un type de filtre et une logique (ET/OU)
    \item Ajoutez plusieurs conditions si nécessaire
\end{enumerate}

\subsection{Types de filtres disponibles}
\begin{itemize}
    \item \textbf{Recherche} : par valeur (texte, nombre, date, booléen)
    \item \textbf{Plage} : intervalle numérique ou temporel
    \item \textbf{Top N} / \textbf{Bottom N}
    \item \textbf{Afficher les vides} / \textbf{Masquer les vides}
    \item \textbf{Tri} : ascendant ou descendant
\end{itemize}

\section{Notebook (rapport)}

\placeholderfigure{Notebook et réglages de rapport}{images/report-notebook.png}{report-notebook}

\subsection{Concept}
Le Notebook regroupe des blocs de contenu (graphique, image, texte). Il apparaît dès qu'un bloc est ajouté et permet d'organiser une analyse puis d'exporter un rapport complet.

\subsection{Ajouter un bloc}
\begin{itemize}
    \item \textbf{Ajouter un graphique}
    \item \textbf{Ajouter une image}
    \item \textbf{Ajouter du texte}
\end{itemize}

\subsection{Éditer et organiser}
\begin{itemize}
    \item Double-clic sur un bloc pour l'éditer
    \item Menu contextuel : déplacer vers le haut/bas, modifier, supprimer
    \item Glisser-déposer pour réordonner les éléments
\end{itemize}

\subsection{Barre d'outils du Notebook}
\begin{itemize}
    \item \textbf{Annuler / Rétablir}
    \item \textbf{Exporter}
\end{itemize}

\subsection{Exporter un rapport}
\begin{itemize}
    \item Formats : \textbf{HTML}, \textbf{PDF}, \textbf{Markdown}
    \item Thèmes : Linen, Ink, Sage, Ember, Harbor, Sandstone
    \item Options : inclure un titre, ouvrir le fichier après export
    \item Prévisualisation intégrée avant l'export
\end{itemize}

\section{Configuration des graphiques}

\placeholderfigure{Paramétrage d'un graphique}{images/chart-settings.png}{chart-settings}

\subsection{Types disponibles}
\begin{itemize}
    \item \textbf{Barres} : comparaison de catégories
    \item \textbf{Colonnes} : tendances et comparaisons verticales
    \item \textbf{Camembert} : proportions
    \item \textbf{Courbe} : évolution continue
    \item \textbf{Nuage de points} : dispersion
    \item \textbf{Aire} : variation cumulée
    \item \textbf{Radar} : comparaison multi-axes
    \item \textbf{Heatmap} : intensité par grille
\end{itemize}

\subsection{Données}
\begin{itemize}
    \item \textbf{Table} : choix de la source (table courante ou autre)
    \item \textbf{Inclure les filtres} : utilise les filtres actifs dans le graphique
    \item \textbf{Pourcentage} : exprime les valeurs en pourcentage
    \item \textbf{Colonne de groupe} (catégorielle) et \textbf{Colonne de valeur} (numérique ou date)
    \item \textbf{Groupe principal} (optionnel) pour hiérarchiser les catégories
\end{itemize}

\subsection{Axes}
\begin{itemize}
    \item Nombre de graduations (ticks)
    \item Échelle : linéaire, logarithmique ou racine carrée
    \item Afficher / masquer l'axe X ou Y
    \item Libellés des axes
\end{itemize}

\subsection{Légende}
\begin{itemize}
    \item Afficher ou masquer la légende
    \item Orientation verticale ou horizontale
\end{itemize}

\subsection{Couleurs}
\begin{itemize}
    \item Sélection de couleurs par série
    \item Prévisualisation via un nuancier
\end{itemize}

\subsection{Options complémentaires}
\begin{itemize}
    \item Titre du graphique
    \item Option \textbf{Empilement} (selon type et disponibilité)
\end{itemize}

\section{Images et texte dans le rapport}

\subsection{Image}
\begin{itemize}
    \item Formats acceptés : PNG, JPG, JPEG, GIF, BMP, WEBP, SVG
    \item Sélectionnez une image via le sélecteur de fichiers
\end{itemize}

\subsection{Texte}
\begin{itemize}
    \item Saisie libre avec retour à la ligne automatique
    \item Idéal pour commentaires, conclusions et annotations
\end{itemize}

\section{Paramètres de l'application}

Les sections de réglages sont repliables pour faciliter la navigation.

\placeholderfigure{Fenêtre des paramètres}{images/settings-dialog.png}{settings-dialog}

\subsection{Général}
\begin{itemize}
    \item Langue de l'interface (10 langues disponibles)
    \item Écran de bienvenue au démarrage
    \item Vérification des mises à jour
    \item Réouverture automatique des tables
    \item Table par défaut (fichier ou URL)
    \item Confirmation à la fermeture de l'application ou d'une table
\end{itemize}

\subsection{Table}
\begin{itemize}
    \item Mode lecture seule
    \item Saisie manuelle (si activée)
    \item Sensibilité de détection de type
    \item Affichage des numéros de ligne
    \item Affichage des valeurs nulles
    \item En-têtes figées (sticky headers)
    \item Masquer les colonnes vides
    \item Afficher les types de colonnes
    \item Redimensionnement automatique des colonnes
    \item Calcul des statistiques de colonnes
\end{itemize}

\subsection{Apparence}
\begin{itemize}
    \item Thème clair, sombre, système ou automatique
    \item Heure de début/fin pour le thème automatique
    \item Activation de FlatLaf
    \item Taille et famille de police
\end{itemize}

\subsection{Raccourcis}
\begin{tabularx}{\textwidth}{@{}l X l@{}}
\toprule
\textbf{Catégorie} & \textbf{Action} & \textbf{Raccourci} \\
\midrule
Table & Ouvrir une table & Ctrl+O \\
Table & Fermer la table & Ctrl+W \\
Affichage & Ouvrir les filtres & Ctrl+F \\
Affichage & Paramètres & Ctrl+, \\
Tri & Tri ascendant & Ctrl+Shift+A \\
Tri & Tri descendant & Ctrl+Shift+D \\
Notebook & Annuler & Ctrl+Z \\
Notebook & Rétablir & Ctrl+Y \\
\bottomrule
\end{tabularx}

\subsection{Avancé}
\begin{itemize}
    \item Journal de débogage
    \item Mode développeur
\end{itemize}

\section{Barre d'état et tâches}
\begin{itemize}
    \item Statut courant (chargement, prêt, erreurs)
    \item Infos sur la table : nom, nombre de lignes, nombre de colonnes
    \item Statistiques rapides sur la colonne sélectionnée
    \item Indicateur de chargement et panneau des tâches en cours (clic sur l'indicateur)
\end{itemize}

\section{Résolution de problèmes}

\subsection{Le fichier ne s'ouvre pas}
\begin{itemize}
    \item Vérifiez l'extension et l'encodage (UTF-8 recommandé)
    \item Assurez-vous que le fichier n'est pas corrompu
    \item Essayez un fichier plus petit pour tester
\end{itemize}

\subsection{Le graphique est vide}
\begin{itemize}
    \item Vérifiez la colonne de valeur (numérique ou date)
    \item Désactivez temporairement les filtres pour tester
    \item Vérifiez la sélection de groupe et de valeur
\end{itemize}

\subsection{Les colonnes disparaissent}
\begin{itemize}
    \item Utilisez \textbf{Afficher toutes les colonnes}
    \item Désactivez \textbf{Masquer les colonnes vides}
\end{itemize}

\section{Support et documentation technique}
\begin{itemize}
    \item \textbf{README.md} : vue d'ensemble du projet
    \item \textbf{DOCUMENTATION.md} : documentation technique complète
    \item \textbf{CONTRIBUTING.md} : guide de contribution
    \item \textbf{MEMBERS.md} : liste des contributeurs
\end{itemize}

\end{document}